%% kapa: exam=1 round=23 year=2012 period=1 form=A subject=economics subject_ko=경제원론 seq=26 number=66
%% source: Q-Net 감정평가사 기출문제 (cst00309) artlSeq=5052578 file=raw/1008/5052578/2012년 제23회 감정평가사 제1차시험 기출문제(1교시 A형).pdf
%% conversion: kordoc 4.24.51 -> tools/gen1.py (auto); review=proofread-sonnet; flags=none
\begin{question}{66}
명목 GDP(국내총생산)와 명목 GNI(국민총소득)에 관한 설명으로 옳은 것을 모두 고른 것은?
\begin{framedtext}
ㄱ. 폐쇄경제에서 명목 GDP는 명목 GNI와 크기가 같다. \par
ㄴ. 한국인이 해외에서 벌어들인 요소소득이 외국인이 한국에서 벌어들인 요소소득보다 더 큰 경우에 명목 GDP가 명목 GNI보다 더 크다. \par
ㄷ. 외국인이 한국에서 벌어들인 근로소득은 한국의 GDP에 포함된다. \par
ㄹ. 한국인이 해외에서 벌어들인 이자수입은 한국의 GDP에 포함된다.
\end{framedtext}
\begin{choices}
  \item ㄱ, ㄴ
  \item ㄱ, ㄷ
  \item ㄴ, ㄷ
  \item ㄴ, ㄹ
  \item ㄷ, ㄹ
\end{choices}
\end{question}
\answer{2}
